\honest{Update Yourself Incrementally}{Eliezer Yudkowsky}{2007}

I open with two slogans from earlier posts: ``Politics is the mind-killer. Debate is war, arguments are soldiers.'' It is tempting to defend a belief against every possible result. You cannot; it is ``mathematically impossible,'' by the rule of my previous post, ``Conservation of Expected Evidence.''

A belief does not need a perfect defence. If a coin lands heads 95\% of the time, it will land tails one time in twenty, and that is normal, as long as there are nineteen heads for each tail. A probabilistic model can take a hit or two and survive, so long as the hits do not keep coming.

Yet it is ``widely believed'' that a true theory can have no failures and a false theory no successes. People hold up one confirming case as if it settled everything, or one contrary case as if it refuted everything. I name no one who argues like this.

Someone will object that you cannot concede a point and still win a debate. My reply: ``Whatever. Rationality is not for winning debates, it is for deciding which side to join.'' I do not dispute the objection's claim about debates. Once you have chosen your side, the work of rationality is already done, well or badly. If choosing the wrong side frightens you at all, you had better integrate all the evidence.

Rationality is a dance in which each step lands exactly where the evidence puts it, down as well as up. If ``an iota or two'' of evidence goes against your belief, ``Just shift downward a little'' and wait for more. \nb{How far one observation should move you depends on what the rival hypothesis predicts, and the post names no rival. If the rival to the 95\% coin is a fair coin, one tail divides the odds for the 95\% coin by ten.} A correct but inexact theory will sometimes meet contrary evidence; only an exact theory is in trouble when it fails once. The trouble with yes-or-no reasoning is that every observation either destroys the theory or does not, so a contrary observation creates dissonance and must be argued away, and incremental progress becomes impossible. On average a correct theory gets more support than countersupport, so you can say without fear that a piece of evidence is gently contrary. ``Think quantitatively.''

Then I restate conservation of expected evidence: ``The weighted mean of your expected posterior probability must equal your prior probability.'' However unlikely a contrary result, the shift it brings must be large enough to balance the expected gain. If you think you already know what the evidence will show, you must already be near certainty, with little room to go up. So it is silly to fear revising downward when you investigate anything at all. \nb{That is the exact form of the rule. It also means that a rare contrary result must move belief a lot, not a little.}

If the contrary evidence keeps coming, you will find you have already given the belief up. ``Yay! Time to celebrate!'' You cannot become stronger by keeping the beliefs you started with.
